\documentclass[11pt,a4paper]{article}
\usepackage[margin=2cm]{geometry}
\usepackage[T1]{fontenc}
\usepackage[utf8]{inputenc}
\usepackage{lmodern}
\usepackage{amsmath,booktabs,tabularx,array,graphicx,float,xcolor,hyperref}
\definecolor{navy}{HTML}{205A83}
\definecolor{rose}{HTML}{AD4166}
\hypersetup{colorlinks=true,urlcolor=navy,linkcolor=navy}
\setlength{\parindent}{0pt}
\setlength{\parskip}{5pt}
\newcommand{\todo}[1]{\par\textcolor{rose}{\textbf{TO COMPLETE:} #1}\par}
\newcommand{\runplots}[2]{%
  \subsection{#1}
  \IfFileExists{results/#2/plots/overview.png}{%
    \begin{figure}[H]\centering
      \includegraphics[width=\linewidth]{results/#2/plots/overview.png}
      \caption{#1: time, performance, estimated effective bandwidth and
      Roofline coordinates. The shaded time range shows the minimum and
      maximum of the five blocks.}
    \end{figure}
  }{\todo{Add #1's overview plot from the selected real run.}}
  \IfFileExists{results/#2/plots/precision_speedup.png}{%
    \begin{figure}[H]\centering
      \includegraphics[width=0.66\linewidth]{results/#2/plots/precision_speedup.png}
      \caption{#1: float32 speedup, defined as $t_{64}/t_{32}$.}
    \end{figure}
  }{\todo{Add #1's precision-speedup plot.}}
}

\title{\textcolor{navy}{Exploring the Roofline Model}}
\author{MohammadJavad Norouzi \and Aland Arif \and
Hamed Zahraei Ramazani \and Wenyuan Huang}
\date{Draft prepared 10 November 2026}

\begin{document}
\maketitle
\textcolor{rose}{\textbf{DRAFT: the four laptop results and final discussion
are still missing. This PDF is not ready to submit.}}

Course: Scientific Programming and High Performance Computing, UPC.
\todo{Add the course repository URL. Update the date and remove the draft
notice only after completing the report.}

\section{Aim}
We compare matrix-vector multiplication, $y=Ax$, and matrix-matrix
multiplication, $C=AB$, using NumPy with \texttt{float32} and \texttt{float64}.
We study how problem size, data reuse and precision affect performance. Each
member runs the same experiment on their own laptop.

\section{Method}
We fill $A$ with ones, $B$ with twos and $x$ with threes. Inputs and output
use the selected precision. We allocate the output before timing and use
\texttt{np.matmul} with its \texttt{out} argument. For these data, each entry
of $y$ should equal $3N$ and each entry of $C$ should equal $2N$. We check the
whole output before and after timing.

The default sizes are $128,256,512,1024,2048,4096$. We perform three warm-up
calls. We then double the calls per block until a block lasts at least
$0.10$ seconds, unless the cap of $1{,}048{,}576$ calls is reached. We measure
five blocks and divide each duration by its call count. We use the median
time per call, while retaining raw times, minima and maxima. Python loop
overhead remains and matters most for small cases.

Allocation, initialization, checking and file output are outside the timer.
NumPy/BLAS may still use internal workspace. We reuse the arrays and do not
flush caches, so this is a warm repeated workload. We request one BLAS thread
and save detected libraries and thread counts. This does not measure the
full multicore peak of the laptop.

\todo{Confirm that these settings describe the selected runs. Record reduced
sizes, interrupted runs, reruns or unverified thread counts.}

\begin{table}[H]\centering\small
\begin{tabularx}{\linewidth}{@{}lXX@{}}\toprule
Member & CPU, RAM and OS & Software, power conditions and selected run \\\midrule
MohammadJavad & TO COMPLETE & TO COMPLETE \\
Aland & TO COMPLETE & TO COMPLETE \\
Hamed & TO COMPLETE & TO COMPLETE \\
Wenyuan & TO COMPLETE & TO COMPLETE \\\bottomrule
\end{tabularx}
\caption{Laptop details. Include Python, NumPy, BLAS and observed threads.}
\end{table}

\section{Theoretical estimates}
A dot product of length $N$ uses $N$ multiplications and $N-1$ additions.
Therefore the exact scalar counts are
\[
F_{\mathrm{mv}}=N(2N-1),\qquad
F_{\mathrm{mm}}=N^2(2N-1).
\]
We use the standard large-$N$ approximations $2N^2$ and $2N^3$ in our
performance calculations. A multiplication and an addition count as two
FLOPs, including when fused into one instruction.

Let $s$ be the bytes per number: $s=4$ for float32 and $s=8$ for float64.
With ideal reuse, matrix-vector multiplication reads $A$ and $x$ once and
writes $y$ once. Matrix-matrix multiplication reads $A$ and $B$ once and
writes $C$ once. Thus
\[
Q_{\mathrm{mv}}=s(N^2+2N),\qquad Q_{\mathrm{mm}}=3sN^2.
\]
We do not read the old output in this model. The resulting estimates are
\[
I_{\mathrm{mv}}=\frac{2N^2}{s(N^2+2N)},\qquad
I_{\mathrm{mm}}=\frac{2N}{3s}.
\]
\begin{table}[H]\centering
\begin{tabular}{@{}lll@{}}\toprule
Operation & float32 intensity & float64 intensity \\\midrule
Matrix-vector, large $N$ & $\approx0.5$ & $\approx0.25$ \\
Matrix-matrix & $N/6$ & $N/12$ \\\bottomrule
\end{tabular}
\caption{Idealized arithmetic intensity in FLOP/byte.}
\end{table}

These estimates omit write-allocate traffic, packing and other internal
transfers. Finite cache may force extra reloads, especially for large
matrix-matrix multiplication. Small warmed cases may obtain inputs from
cache. The estimated bytes are not measured RAM traffic.

For median time $t$, we calculate
\[
P=\frac{F}{t\,10^9}\ \text{GFLOP/s},\qquad
B_{\mathrm{est}}=\frac{Q}{t\,10^9}\ \text{GB/s}.
\]
GB is decimal. Since $B_{\mathrm{est}}=P/I$, these quantities are not
independent measurements of computation and physical memory bandwidth.

\section{Roofline interpretation}
The Roofline model is
\[
P\leq\min\left(P_{\mathrm{peak}},B_{\mathrm{memory}}I\right).
\]
The bandwidth and intensity must refer to the same memory level. The ridge
is at $I=P_{\mathrm{peak}}/B_{\mathrm{memory}}$. Low intensity makes memory
supply a likely limit; high intensity can make the compute ceiling lower
if data are reused effectively. Actual performance can be below both ceilings.

We do not independently measure hardware ceilings in the basic experiment.
Our default plots show measured GFLOP/s versus estimated intensity without
claimed hardware-limit lines. This supports a qualitative interpretation,
but does not prove the limiting resource. We do not treat the largest value
from our benchmarks as a hardware ceiling.
\todo{If independently sourced ceilings were added, give their sources,
memory level, precision and matching thread setting. Otherwise remove this
instruction and keep the preceding limitation.}

\section{Experimental results}
\todo{Select one real run per laptop. The standard folders below are
mohammadjavad, aland, hamed and wenyuan. Change the folder arguments in the
LaTeX source if a rerun was selected.}

\runplots{MohammadJavad}{mohammadjavad}
\runplots{Aland}{aland}
\runplots{Hamed}{hamed}
\runplots{Wenyuan}{wenyuan}

\subsection{Numerical comparison}
\todo{Choose the largest common $N$ and fill in the table from the CSV files.}
\begin{table}[H]\centering\small
\begin{tabular}{@{}lrrrrr@{}}\toprule
Member & $N$ & MV $t_{64}/t_{32}$ & MM $t_{64}/t_{32}$ & MV FP64 & MM FP64 \\\midrule
MohammadJavad & -- & -- & -- & -- & -- \\
Aland & -- & -- & -- & -- & -- \\
Hamed & -- & -- & -- & -- & -- \\
Wenyuan & -- & -- & -- & -- & -- \\\bottomrule
\end{tabular}
\caption{Precision speedup and FP64 performance. The last two columns use
GFLOP/s. A speedup above one means float32 was faster.}
\end{table}

\section{Discussion}
\subsection{How does time scale with problem size?}
The work grows as $N^2$ for matrix-vector and $N^3$ for matrix-matrix.
Doubling $N$ increases the work by about four and eight respectively.
This is the expected time scaling only while achieved performance stays
approximately constant.
\todo{Give actual ratios for both operations and precisions, naming the
sizes. Discuss overhead, cache transitions and changing GFLOP/s using evidence.}

\subsection{How do the operations compare?}
Matrix-vector has low, nearly constant ideal intensity. Matrix-matrix can
reuse an input entry for many outputs, giving intensity that grows with $N$.
Optimized BLAS uses blocking to improve reuse.
\todo{Compare measured GFLOP/s at the same $N$. Also compare times: higher
GFLOP/s does not mean less time when the operation does more total work.}

\subsection{What changes with float32?}
Float32 halves array bytes and doubles the ideal intensity. It may also
increase vector arithmetic throughput. A twofold speedup is not guaranteed.
\todo{Report $t_{64}/t_{32}$ for both operations. Discuss differences between
sizes and laptops and any cases in which float32 was not faster.}

\subsection{Were there unexpected results?}
\todo{Identify unusual points or large timing spreads, or state that no
substantial irregularity was found. Possible causes include background work,
clock and power behavior, cache capacity, BLAS choices and timing overhead.
Distinguish measured evidence from possible explanations.}
We keep the four laptops separate. Hardware and software vary together,
so this comparison does not isolate one hardware feature by itself.

\section{Conclusion}
\todo{Write a short conclusion from the real measurements: the operation
difference, precision effect and limits of the estimated-byte interpretation.}

\section{Contributions}
\todo{Use the four contribution notes and PR links. Record sample-code review,
changes, running, checking plots and report work accurately.}

\section*{References}
\begin{enumerate}
\item Course handout, \emph{Homework: Exploring the Roofline Model}.
\item NumPy, \texttt{matmul}:\\
\url{https://numpy.org/doc/stable/reference/generated/numpy.matmul.html}
\item threadpoolctl: \url{https://github.com/joblib/threadpoolctl}
\item NERSC, Roofline Performance Model:\\
\url{https://docs.nersc.gov/tools/performance/roofline/}
\end{enumerate}
\end{document}
